% In this file you should put the actual content of the blueprint.
% It will be used both by the web and the print version.
% It should *not* include the \begin{document}
%
% If you want to split the blueprint content into several files then
% the current file can be a simple sequence of \input. Otherwise It
% can start with a \section or \chapter for instance.

\chapter{The domain of complex metrics (KS Section 2)}

This chapter is the blueprint for Section~2 of Kontsevich--Segal,
\emph{Wick rotation and the positivity of energy in quantum field theory}
(arXiv:2105.10161), formalized in \texttt{KontsevichSegal/ComplexMetrics/}.
The paper introduces, for a $d$-dimensional real vector space $V$, an open
domain $Q_{\mathbb C}(V)$ of \emph{allowable} complex-valued quadratic forms
(``complex metrics''), generalizing positive-definite inner products. The
Shilov boundary of this domain contains the real Lorentzian metrics and no
other nondegenerate real metrics.

\section{Allowable complex metrics}

\begin{definition}[Allowable complex metric, Hodge-star form; KS Definition 2.1]
  \label{def:allowable-hodge}
  \notready
  On a $d$-dimensional real vector space $V$, a quadratic form
  $g \colon V \to \mathbb C$ is an \emph{allowable complex metric} if, for
  all degrees $p \geq 0$, the real part of the quadratic form
  \[
    \Lambda^p(V^*) \longrightarrow |\Lambda^d(V^*)|_{\mathbb C}, \qquad
    \alpha \mapsto \alpha \wedge *_g\, \alpha
  \]
  is positive-definite. Here $*_g$ is the Hodge star operator of the complex
  metric $g$, and $|\Lambda^d(V^*)|_{\mathbb C}$ is the complexified twisted
  determinant line, with orientation-reversing automorphisms of $V$ acting
  antilinearly.

  The Hodge star operator for complex metrics, and the induced quadratic
  forms on the exterior powers $\Lambda^p(V^*)$, are not available in
  Mathlib.
\end{definition}

\begin{definition}[Allowable complex metric, working definition: the angle condition]
  \label{def:allowable}
  \lean{AngleCondition, AllowableComplexMetric, QC}
  \leanok
  A tuple of complex numbers $(\lambda_1, \dots, \lambda_d)$ satisfies the
  \emph{angle condition} if every $\lambda_i$ is nonzero, no $\lambda_i$ lies
  on the non-positive real axis, and
  \[
    |\arg(\lambda_1)| + |\arg(\lambda_2)| + \dots + |\arg(\lambda_d)|
    \;<\; \pi.
    \tag{4}
  \]
  An \emph{allowable complex metric} on a finite-dimensional real vector
  space $V$ is a symmetric nondegenerate $\mathbb R$-bilinear form
  $g \colon V \times V \to \mathbb C$ for which there exist a basis
  $(y_i)$ of coordinates on $V$ and coefficients $(\lambda_i)$ satisfying
  the angle condition with
  \[
    g \;=\; \lambda_1 y_1^2 + \lambda_2 y_2^2 + \dots + \lambda_d y_d^2 .
  \]
  The collection of allowable complex metrics on $V$ is denoted
  $Q_{\mathbb C}(V)$; in the paper it is an open subset of the complex
  vector space $S^2(V^*_{\mathbb C})$. This diagonal characterization serves
  as the definition of allowability throughout the formalization, since
  Definition~\ref{def:allowable-hodge} requires the Hodge star operator.
\end{definition}

\begin{theorem}[Equivalence of the two definitions; KS Theorem 2.2]
  \label{thm:angle-equiv}
  \lean{defn_2_1_equiv_angle_condition}
  \notready
  \uses{def:allowable-hodge, def:allowable}
  A quadratic form $g \colon V \to \mathbb C$ satisfies
  Definition~\ref{def:allowable-hodge} if and only if it can be written, in
  some basis of $V$, as $\sum_i \lambda_i y_i^2$ with nonzero coefficients
  $\lambda_i$, not on the negative real axis, satisfying the angle
  condition (4) of Definition~\ref{def:allowable}.

  The biconditional cannot yet be stated in Lean: the Hodge-star side
  (Definition~\ref{def:allowable-hodge}) is not expressible, pending the
  Hodge star operator and induced forms on exterior powers in Mathlib.
\end{theorem}

\begin{proof}
  \notready
  \uses{def:allowable-hodge, def:allowable}
  (Paper, p.~10.) A form satisfying Definition~\ref{def:allowable-hodge} with
  $p = 1$ has $\operatorname{Re}\bigl((\det g)^{-1/2} g\bigr)$
  positive-definite, which allows simultaneous diagonalization of the real
  and imaginary parts of $g$ over $\mathbb R$. In a diagonalizing basis the
  form $\alpha \mapsto \alpha \wedge *_g \alpha$ on $\Lambda^p(V^*)$ is
  diagonal with values
  $(\lambda_1 \cdots \lambda_d)^{1/2} \prod_{i \in S} \lambda_i^{-1}$ on the
  basis $p$-forms $e_S^*$; positivity of all their real parts for every
  subset $S$ is precisely condition (4).
\end{proof}

\begin{lemma}[Values form a convex cone; KS p.~9]
  \label{lem:convex-cone}
  \lean{not_neg_real_axis}
  \leanok
  \uses{def:allowable}
  For $g \in Q_{\mathbb C}(V)$ one has
  $\max_i \arg \lambda_i - \min_i \arg \lambda_i < \pi$, so as $v$ runs
  through $V$ the values $g(v)$ form a closed convex cone in $\mathbb C$
  disjoint from the open negative real axis. In particular $g(v)$ is never
  real and negative.
\end{lemma}

\begin{proof}
  \leanok
  \uses{def:allowable}
  From (4), if $\arg \lambda_j > 0 > \arg \lambda_k$ then
  $\arg \lambda_j - \arg \lambda_k \leq \sum_i |\arg \lambda_i| < \pi$;
  hence all values $\lambda_i y_i^2$ lie in a closed convex cone of opening
  angle less than $\pi$, which contains no two opposite rays and avoids the
  open negative real axis. Nonnegative combinations stay in the cone.
\end{proof}

\begin{lemma}[The volume element; KS p.~7]
  \label{lem:volume-element}
  \lean{volume_element_positive}
  \leanok
  \uses{def:allowable}
  For $g \in Q_{\mathbb C}(V)$ the determinant $\det g$ (of the Gram matrix
  in any basis; it is invariantly defined up to multiplication by a positive
  real number) is not real and negative, and the square root
  $(\det g)^{1/2}$ can be chosen with strictly positive real part. This
  fixes the holomorphic volume element
  $\mathrm{vol}_g = (\det g)^{1/2} |dy^1 \cdots dy^d|$.
\end{lemma}

\begin{proof}
  \leanok
  \uses{def:allowable}
  In a diagonalizing basis $\det g = \lambda_1 \cdots \lambda_d$, whose
  argument is $\sum_i \arg \lambda_i \in (-\pi, \pi)$ by (4); a base change
  multiplies the Gram determinant by $(\det P)^2 > 0$. Since
  $\arg \det g \in (-\pi,\pi)$, the principal square root has
  $|\arg| < \pi/2$, i.e.\ positive real part.
\end{proof}

\begin{proposition}[Configurations in the holomorphic envelope; KS Proposition 2.3]
  \label{prop:conf-envelope}
  \notready
  \uses{def:allowable}
  If a complex metric on a $d$-dimensional real vector space $V$ satisfies
  condition (4), then the configuration space $\mathrm{Conf}_k(V)$ of
  $k$-tuples of distinct points of $V$ is contained in the holomorphic
  envelope of the Wightman permuted extended tube
  $\mathcal U_k(V) \subset (V_{\mathbb C})^k$.

  Formalization requires the Wightman permuted extended tube, holomorphic
  envelopes of domains in $(\mathbb M_{\mathbb C})^k$, and Bochner's tube
  theorem, none of which are available; deferred to the
  \texttt{WickRotation} phase.
\end{proposition}

\begin{proof}
  \uses{lem:l-shape}
  (Paper, appendix to Section 2.) Identify $V$ with
  $\exp(\mathrm i \Theta/2)(\mathbb E) \subset \mathbb M_{\mathbb C}$ where
  $\Theta$ lies in the generalized octahedron $\Pi_0$ cut out by (4). One
  shows $\exp(\mathrm i \Theta/2)(\mathbf x) \in \mathcal U_k$ when
  $\mathrm{Re}(\Theta)$ lies on a coordinate axis ($\Pi_{00}$), then extends
  to the convex hull $\Pi_0$ of $\Pi_{00}$ by Bochner's tube theorem; the
  reduction of the $d = 2$ case to the standard tube theorem is
  Lemma~\ref{lem:l-shape}.
\end{proof}

\section{Structure of the domain $Q_{\mathbb C}(V)$}

\begin{definition}[The fiber $\Pi(V)$: trace-norm ball; KS p.~11]
  \label{def:trace-norm-fiber}
  \lean{TraceNormLtOne}
  \leanok
  For a positive-definite inner product $g_0$ on $V$, let $\Pi(V, g_0)$
  denote the convex open set of $g_0$-self-adjoint operators
  $\Theta \colon V \to V$ with trace-norm
  $\|\Theta\|_1 = \sum_i |\theta_i| < 1$, where $\theta_i$ are the
  eigenvalues; that is, the interior of the convex hull of the rank-one
  orthogonal projections. Since $V$ carries no inner-product structure in
  the formalization, self-adjointness with respect to \emph{some}
  positive-definite $g_0$ is encoded equivalently as diagonalizability over
  $\mathbb R$: an eigenbasis with real eigenvalues $\theta$ satisfying
  $\sum_i |\theta_i| < 1$.
\end{definition}

\begin{proposition}[Parametrization and contractibility; KS Proposition 2.4]
  \label{prop:parametrization}
  \lean{QC_parametrization, QC_contractible}
  \notready
  \uses{def:allowable, def:trace-norm-fiber}
  $Q_{\mathbb C}(V)$ is a fibre bundle over the space of positive-definite
  inner products on $V$ whose fibre at $g_0$ is $\Pi(V, g_0)$. Equivalently,
  choosing a reference inner product,
  \[
    Q_{\mathbb C}(V) \;\cong\; \mathrm{GL}(V) \times_{\mathrm{O}(V)} \Pi(V).
  \]
  In particular, $Q_{\mathbb C}(V)$ is contractible.

  Stating this in Lean requires $\mathrm{GL}(V)/\mathrm{O}(V)$ actions with
  associated-bundle machinery, an inner-product-space structure compatible
  with the algebraic setup of the formalization, and a topology on
  $Q_{\mathbb C}(V)$, which is needed even to state contractibility; none
  are available in Mathlib.
\end{proposition}

\begin{proof}
  \notready
  \uses{def:allowable}
  (Paper, pp.~10--11.) The diagonalization of
  Definition~\ref{def:allowable} presents $g$ as a sequence of angles
  $\theta_1 \geq \dots \geq \theta_m$ in $(-\pi,\pi)$ and a decomposition
  $V = V_1 \oplus \dots \oplus V_m$ with
  $\sum_k \dim V_k \cdot |\theta_k| < \pi$, on $V_k$ the form being
  $e^{\mathrm i \theta_k}$ times a positive-definite real form. The subspace
  $P = \bigoplus e^{-\mathrm i \theta_k/2} V_k \subset V_{\mathbb C}$ is
  canonical, carries a real positive-definite form $g_0$, and
  $V = \exp(\mathrm i \pi \Theta/2)(P)$ where $\Theta$ is $g_0$-self-adjoint
  with $\|\Theta\|_1 < 1$; condition (4) is exactly the trace-norm bound.
\end{proof}

\begin{definition}[Restriction of the form to a subspace]
  \label{def:restrict-form}
  \lean{AllowableComplexMetric.restrict, AllowableComplexMetric.restrict_symmetric}
  \leanok
  \uses{def:allowable}
  For $g \in Q_{\mathbb C}(V)$ and a vector subspace $W \leq V$, the
  restriction $g|_W \colon W \times W \to \mathbb C$,
  $(w_1, w_2) \mapsto g(w_1, w_2)$, is a symmetric $\mathbb R$-bilinear
  form on $W$.
\end{definition}

\begin{proposition}[Restriction preserves allowability; KS Proposition 2.5]
  \label{prop:restriction}
  \lean{restrict_allowable}
  \leanok
  \uses{def:allowable, def:restrict-form}
  If $g \in Q_{\mathbb C}(V)$ and $W$ is any vector subspace of $V$, then
  $g|_W$ belongs to $Q_{\mathbb C}(W)$.
\end{proposition}

\begin{proof}
  \uses{def:allowable}
  (Paper, p.~13.) For $g \in Q_{\mathbb C}(V)$ the map
  $v \mapsto \arg g(v)$ is smooth on $\mathbb P(V)$ with critical values the
  angles $\theta_1 \geq \dots \geq \theta_d$, characterized by a min--max
  over subspaces. It suffices to treat $W$ of codimension one; then the
  critical values $\theta'_k$ of $\arg(g|_W)$ \emph{interleave}:
  $\theta_k \geq \theta'_k \geq \theta_{k+1}$, whence
  $\sum_k |\theta'_k| \leq \sum_k |\theta_k| < \pi$.
\end{proof}

\begin{proposition}[Real allowable subspaces of a complex quadratic space; KS Proposition 2.6]
  \label{prop:real-subspaces}
  \notready
  \uses{def:allowable, def:trace-norm-fiber}
  Let $Z$ be a $d$-dimensional complex vector space with a nondegenerate
  quadratic form $g$, and let $\mathcal R(Z)$ be the space of
  $d$-dimensional real subspaces $A \subset Z$ such that $g|_A$ belongs to
  $Q_{\mathbb C}(A)$ (an open subset of the real Grassmannian). Then
  $\mathcal R(Z)$ is contractible, and
  \[
    \mathcal R(Z) \;\cong\;
    \mathrm{O}_{\mathbb C}(Z) \times_{\mathrm{O}(Z_{\mathbb R})}
    \Pi(Z_{\mathbb R}).
  \]

  Formalization requires the Grassmannian of real subspaces,
  $\mathrm{O}_{\mathbb C}(Z)$-actions, and the same associated-bundle
  machinery as Proposition~\ref{prop:parametrization}. The result is needed
  for the Wick rotation of field operators in Section~5 of the paper and is
  deferred to \texttt{WickRotation/FieldOperators}.
\end{proposition}

\begin{proof}
  \uses{prop:parametrization}
  (Paper, p.~13.) The space $\mathcal R(V; Z)$ of $\mathbb R$-linear
  embeddings $f \colon V \to Z$ with $f^*(g)$ allowable is simultaneously a
  principal $\mathrm{GL}(V)$-bundle over $\mathcal R(Z)$ and a principal
  $\mathrm{O}_{\mathbb C}(Z)$-bundle over $Q_{\mathbb C}(V)$. The
  contractibility of $Q_{\mathbb C}(V)$
  (Proposition~\ref{prop:parametrization}) and the homotopy equivalence of
  both structure groups with the compact $\mathrm O_d$ yield the
  contractibility of $\mathcal R(Z)$.
\end{proof}

\begin{proposition}[Domain of holomorphy; KS Proposition 2.7]
  \label{prop:domain-holomorphy}
  \lean{QC_domain_of_holomorphy}
  \notready
  \uses{def:allowable}
  The domain $Q_{\mathbb C}(V)$ is holomorphically convex, i.e.\ a domain of
  holomorphy in $S^2(V^*_{\mathbb C})$.

  Stating this in Lean requires a complex-manifold structure on the space of
  quadratic forms, the notion of domain of holomorphy (Stein theory), and
  Siegel domains; none of these are in Mathlib.
\end{proposition}

\begin{proof}
  \notready
  \uses{def:allowable-hodge}
  (Paper, p.~14.) The Siegel domain $\mathcal U(V)$ of complex inner
  products with positive-definite real part is a domain of holomorphy; so is
  the product $\prod_{0 \leq p \leq d/2} \mathcal U(\Lambda^p(V))$.
  $Q_{\mathbb C}(V)$ is the intersection of this product domain with the
  affine subvariety given by the natural embedding of
  $S^2(V^*_{\mathbb C})$ (whose components are the exterior-power forms of
  Definition~\ref{def:allowable-hodge}), hence itself a domain of
  holomorphy.
\end{proof}

\begin{lemma}[Holomorphic extension from an L-shaped tube; KS Lemma 2.8]
  \label{lem:l-shape}
  \notready
  Let $L$ be the L-shaped subset
  $(\{0\} \times [0,1)) \cup ([0,1) \times \{0\})$ of the quadrant
  $(\mathbb R_+)^2$. Then any holomorphic function $F$ defined in a
  neighbourhood of $L \times \mathrm i \mathbb R^2 \subset \mathbb C^2$
  extends holomorphically to $P \times \mathrm i \mathbb R^2$, where $P$ is
  the intersection of $(\mathbb R_+)^2$ with a neighbourhood of $L$ in
  $\mathbb R^2$.

  Formalization requires several-complex-variables extension theory
  (continuity-principle and analytic-disc arguments, Bochner's tube
  theorem), not in Mathlib; deferred together with
  Proposition~\ref{prop:conf-envelope}.
\end{lemma}

\section{The Shilov boundary: Lorentzian metrics}

\begin{definition}[Lorentzian signature]
  \label{def:lorentzian}
  \lean{IsLorentzian}
  \leanok
  A real symmetric bilinear form $\varphi$ on $V$ is \emph{Lorentzian} (of
  signature $(d-1, 1)$) if in some basis it diagonalizes as
  $\sum_i s_i x_i y_i$ with each $s_i \in \{+1, -1\}$ and exactly one
  $s_i = -1$.
\end{definition}

\begin{proposition}[Lorentzian metrics lie on the boundary; KS p.~9]
  \label{prop:lorentzian-boundary}
  \lean{lorentzian_on_boundary}
  \leanok
  \uses{def:allowable, def:lorentzian}
  A real Lorentzian metric $\varphi$ on $V$ lies on the boundary of
  $Q_{\mathbb C}(V)$: it is a limit of allowable complex metrics
  (entry-wise on a fixed basis, i.e.\ in the canonical topology on the
  finite-dimensional space of forms), but is not itself allowable.
\end{proposition}

\begin{proof}
  \leanok
  \uses{def:allowable, lem:convex-cone}
  Diagonalize $\varphi = \sum_i s_i y_i^2$ with exactly one $s_i = -1$ and
  perturb the negative coefficient to $e^{\mathrm i (\pi - \delta)}$: the
  angle sum is $\pi - \delta < \pi$, so the perturbation is allowable and
  converges to $\varphi$ as $\delta \downarrow 0$. That $\varphi$ itself is
  not allowable follows from Lemma~\ref{lem:convex-cone}: $\varphi$ takes a
  real negative value on the timelike basis vector.
\end{proof}

\begin{proposition}[Only Lorentzian metrics lie on the boundary; KS p.~9]
  \label{prop:only-lorentzian}
  \lean{only_lorentzian_on_boundary}
  \leanok
  \uses{def:allowable, def:lorentzian}
  If a real nondegenerate bilinear form $\varphi$ on $V$ is a limit of
  allowable complex metrics but is not positive-definite (the
  positive-definite forms being the real points of the open domain
  $Q_{\mathbb C}(V)$ itself), then $\varphi$ is Lorentzian. Equivalently,
  among real nondegenerate metrics, those of signature $(d-1,1)$, and no
  others, lie on the boundary of $Q_{\mathbb C}(V)$.
\end{proposition}

\begin{proof}
  \uses{def:allowable, prop:restriction}
  For a real metric each $|\arg \lambda_i|$ is $0$ or $\pi$, and inequality
  (4) shows at most one $|\arg \lambda_i|$ can become $\pi$ in the limit: a
  two-dimensional $\varphi$-negative-definite subspace would, after
  restricting the approximating allowable metrics to it
  (Proposition~\ref{prop:restriction}), force the angle sum toward $2\pi$.
  Nondegeneracy and failure of positive-definiteness then force exactly one
  negative eigenvalue.
\end{proof}

\begin{theorem}[Two copies of the Lorentzian metrics on the Shilov boundary; KS p.~9]
  \label{thm:two-copies}
  \lean{two_copies_on_boundary}
  \notready
  \uses{def:allowable, def:lorentzian, prop:domain-holomorphy}
  The Shilov boundary of $Q_{\mathbb C}(V)$, regarded as a bounded domain in
  an affine variety as in the proof of
  Proposition~\ref{prop:domain-holomorphy}, contains \emph{two} disjoint
  copies of the space of Lorentzian metrics on $V$, for an eigenvalue
  $\lambda$ can approach the negative real axis either from above or from
  below. The two copies are interchanged by the complex-conjugation map on
  $Q_{\mathbb C}(V)$. With orientation-reversing elements of
  $\mathrm{GL}(V)$ acting antilinearly on the orientation line, the
  nondegenerate points of the Shilov boundary are the \emph{time-oriented}
  Lorentzian metrics.

  The faithful statement is not yet expressible in Lean. It requires:
  (1) induced forms on exterior powers and the Hodge star, for the
  affine-variety embedding
  $Q_{\mathbb C}(V) \to \prod_{0 \leq p \leq d/2} \mathcal U(\Lambda^p V)$
  from the proof of Proposition~\ref{prop:domain-holomorphy}, whose $p = 0$
  factor $(\det g)^{1/2}$ carries the branch that separates the two copies
  (in the naive closure inside $S^2(V^*_{\mathbb C})$ the copies coincide);
  (2) Siegel domains and the Cayley transform, giving the bounded
  realization; (3) the Shilov boundary notion itself, the smallest compact
  subset of the closure on which every holomorphic function attains its
  supremum.
\end{theorem}

\section{The two-dimensional case}

\begin{proposition}[$Q_{\mathbb C}(V)$ is a polydisc when $\dim V = 2$; KS p.~15]
  \label{prop:two-dim-polydisc}
  \lean{QC_two_dim_polydisc}
  \notready
  \uses{def:allowable}
  When $\dim V = 2$, the conformal class $(\det g)^{-1/2} g$ decouples from
  the volume element, and a nondegenerate complex metric is determined up to
  scale by its two null directions in $\mathbb P(V_{\mathbb C})$;
  allowability puts one in each open hemisphere separated by the equator
  $\mathbb P(V)$. The domain $Q_{\mathbb C}(V)$ is thus a $3$-dimensional
  polydisc: one disc for each of the two complex structures (null
  directions) and one for the volume element.

  Stating this in Lean requires the projective geometry of
  $\mathbb P(V_{\mathbb C})$ and the conformal/volume decomposition of
  quadratic forms.
\end{proposition}
